%! The Language of Variation — Unit 1, Lesson 1.1 — Student Packet Cover Sheet
\documentclass[10pt]{article}
\usepackage{apstats-article}
\usepackage{apstats-boxes}
\usepackage{ltablex}
\keepXColumns

\begin{document}

% Full-bleed navy banner
\begin{tikzpicture}[remember picture, overlay]
  \fill[navy] (current page.north west)
    rectangle ([yshift=-0.9in]current page.north east);
\end{tikzpicture}
\vspace{-0.8in}

\begin{center}
  {\color{white}\textbf{\LARGE Statistics}}\\[4pt]
  {\color{skymid}\large\textbf{Unit 1}}\\[6pt]
  {\color{white}\textbf{\large Lesson 1.1 \quad The Language of Variation}}
\end{center}

\vspace{0.22in}
\namedateperiod
\vspace{0.18in}

\boxguard[14]
\begin{learningtargetbox}
\small
% GRADUAL RELEASE: the vocabulary is taught in the guided notes, so the targets name it
% plainly. The AP Learning Objectives (VAR-1.B, VAR-1.C) stay in the teacher-facing plan.
\begin{enumerate}[leftmargin=1.5em, itemsep=5pt]
  \item \ldots{} classify each variable in a data set as \textbf{categorical} or
        \textbf{quantitative}.
  \item \ldots{} use the \textbf{amount test} --- does the value record how much or how many?
        --- instead of asking whether it is written in digits.
  \item \ldots{} explain why averaging a ZIP code or a permit number is meaningless.
  \item \ldots{} explain what changes, and what is lost, when measured amounts are replaced by
        group labels.
\end{enumerate}
\end{learningtargetbox}

\vspace{0.14in}

\boxguard[14]
\begin{tocbox}
\small
\renewcommand{\arraystretch}{1.5}
\begin{tabularx}{\linewidth}{@{} c l X r @{}}
\rowcolor{sky}
\textbf{\#} & \textbf{Component} & \textbf{Description} & \textbf{Score} \\
\midrule
1 & Warm-Up & An instrument rental sheet, and four lists --- which ones can you add up?
                 & \blank{1.2cm} \\
2 & Guided Notes \& Practice & The two kinds of variable, the amount test, and the columns
                 written in digits that fool people --- then practice together and practice alone
                 & \blank{1.2cm} \\
3 & Group Activity & \emph{The Trailhead Register} --- two problems with your group, then the
                 debrief we fill in together & \blank{1.2cm} \\
% AP Practice is assigned but NOT scored: its score cell prints NA, never a \blank{}.
% Row 5 is the lesson's GRADED practice; when DeltaMath carries this topic instead, that
% assignment replaces row 5 and its score cell is struck through.
4 & AP Practice & The trailhead register again: AP-style multiple choice and one free-response
                 set --- practice, not a grade & \textbf{NA} \\
5 & Homework & A coffee-shop order log, a pet survey, a hospital chart, and one
                 multiple-choice item --- \textbf{scored, due next class} & \blank{1.2cm} \\
\midrule
  & \multicolumn{2}{r}{\textbf{Total}} & \blank{1.2cm} \\
\end{tabularx}
\end{tocbox}

\vspace{0.14in}

\boxguard[8]
\begin{remindbox}
\small
Today runs in three moves: \textbf{I show you}, \textbf{we do one together}, \textbf{you do
it with your group}. The guided notes give you the two kinds of variable and the test that
separates them; the trailhead register is where you use it. \textbf{Some of the columns are
traps} --- trust what a column \emph{records}, not what it looks like.
\end{remindbox}

\end{document}
